\section{Group Theory}

\begin{question}[subtitle = January 2026 Problem 5,ID=Jan26_5]
Let $G$ be a nontrivial finite group. A normal subgroup $N \lnormaleq G$ is \emph{minimal} if $N \neq \set{e}$ and it contains no proper subgroups normal in $G$. For example, $G$ itself is a minimal normal subgroup if and only if it is simple.
\begin{enumerate}
    \item Let $N_1, N_2 \lnormaleq G$ be distinct minimal normal subgroups of $G$. Show that the subgroup $\genset{N_1, N_2}$ generated by $N_1$ and $N_2$ is isomorphic to $N_1 \times N_2$.
    \item Let $N_1, \dots, N_r \lnormaleq G$ be minimal normal subgroups, $r \geq 1$. Put $H = \genset{N_1,\dots, N_r}$. Prove that
    \[H \cong \prod_{i \in S} N_i\]
    for some subset $S \subset \set{1,\dots,r}$.
    \item Let $N \lnormaleq G$ be a minimal normal subgroup, and let $H \lnormaleq N$ be any normal subgroup, $H \neq \set{e}$. Show that $N$ is generated by the conjugates of $H$:
    \[N = \genset{\bigcup_{g \in G} gHg\inv}.\]
    \item Show that for any minimal normal subgroup $N \lnormaleq G$, there exists a simple group $S$ and $k > 0$ such that $N \cong S^k$.
\end{enumerate}
\end{question}
\begin{solution}
\begin{enumerate}
    \item Since the intersection of two normal subgroups is again normal, $N_1 \cap N_2$ is a normal subgroup of $G$ properly contained in both, hence $N_1 \cap N_2 = \set{e}$ by the minimality assumption. Thus, $\genset{N_1, N_2} \cong N_1 \times N_2$ by the internal direct product test.
    \item We proceed inductively from the $r = 2$ case above. Note that a subgroup generated by normal subgroups is again normal. So, if we have that
    \[H' \coloneqq \genset{N_1,\dots,N_{r-1}} = \prod_{i \in S'} N_i\]
    for $S' \subset \set{1,\dots,r-1}$, then $H' \cap N_r$ is either trivial or all of $N_r$. In the former case, $\genset{H',N_r} = H' \times N_r$, and in the latter case $\genset{H',N_r} = H'$.
    \item Note that the subgroup generated by the conjugates of $H$ is preserved by conjugation in $G$, by construction, hence is a nontrivial normal subgroup of $G$. On the other hand, since $gHg\inv \subset gNg\inv = N$, the subgroup generated by the conjugates of $H$ is contained in $N$. Since $N$ is minimal, we must have equality. (We did not actually use that $H$ was normal in $N$, but the hypothesis was probably included to be suggestive for the next part.)
    \item Let $S$ be a minimal normal subgroup of $N$. By part (c), the $G$-conjugates of $S$ generate $N$. Since each conjugate of $S$ is again a minimal normal subgroup of $N$, by part (b) we have that $N$ is the internal direct product of conjugates of $S$. Since each conjugate of $S$ is isomorphic to $S$, this tells us that $N \cong S^k$ for some $k$.
    
    Finally, we will show that $S$ is simple. Suppose that $H \lnormal S$, then $H$ is normal in a direct factor of $N$, hence $H \lnormal N$. But we supposed that $S$ was a minimal normal of $N$, so we must have either $H = S$ or $H = \set{e}$. Thus, $S$ is simple. (I think it's somewhat funny that we show $S$ is simple by first showing that $N \cong S^k$. I tried to think how to show $S$ is simple directly, and it doesn't seem easy.)
\end{enumerate}
\end{solution}

%
\begin{question}[subtitle=August 2025 Problem 3,ID=Aug25_3]
Let $G$ be a finite group.
\begin{enumerate}
    \item Suppose $H \subset G$ is a normal subgroup such that the quotient group $G/H$ is cyclic of order $n$. Show that for every divisor $d$ of $n$, there exists a normal subgroup $K \subset G$ such that the quotient group $G/K$ is cyclic of order $d$.
    \item Suppose $H_1, H_2 \subset G$ are two normal subgroups such that the quotients are cyclic:
    \[G/H_1 \simeq \Z/n_1\Z \qquad G/H_2 \simeq \Z/n_2\Z.\]
    Assume further that $n_1$ and $n_2$ are coprime. Show that there exists a normal subgroup $H \subset G$ such that
    \[G/H \simeq \Z/(n_1 n_2 \Z).\]
\end{enumerate}
\end{question}
\begin{solution}
\begin{enumerate}
    \item By the lattice isomorphism theorem, normal subgroups of $G/H$ correspond to normal subgroups of $G$ which contain $H$. Since $G/H$ is cyclic, for every $d$ dividing $n$, it contains a normal subgroup $K'$ of index $d$, and the quotient will be cyclic. This normal subgroup corresponds to a normal subgroup $K$ of $G$, and by the third isomorphism theorem(!) $G/K \cong (G/H)/K'$ is cyclic of order $d$.
    \item (Solution courtesy of Caroline Nunn.) Let $\pi_1\colon G \to G/H_1$ and $\pi_2\colon G \to G/H_2$ be the quotient maps. We have a map $D\colon G \to G/H_1 \times G/H_2$ given by $D(g) = (\pi_1(g),\pi_2(g))$. We have $G/H_1 \times G/H_2 \cong (\Z/n_1\Z) \times (\Z/n_2\Z) \cong \Z/(n_1n_2 \Z)$ by the Chinese Remainder Theorem. So, if we can show that $D$ is surjective, then we will be done, since $\ker D$ will be our desired subgroup.
    
    We will show that this map is surjective by showing that we can find elements $g_1,g_2 \in G$ so that $D(g_1) = (\mathrm{generator},e)$ and $D(g_2) = (e,\mathrm{generator})$. Let $x \in G$ be any element so that $\pi_1(x)$ generates $G/H_1$. Since $n_1$ and $n_2$ are coprime, $\pi_1(x^{n_2})$ is also a generator of $G/H_1$, but now $\pi_2(x^{n_2}) = \pi(x)^{n_2} = e$. Similarly, if $\pi_2(y)$ is a generator for $G/H_2$, so is $\pi_2(y^{n_1})$, but $\pi_1(y^{n_1}) = e$. Thus, the choice $g_1 = x^{n_2}$, $g_2 = y^{n_1}$ does what we wanted.
\end{enumerate}
\end{solution}

%
\begin{question}[subtitle=January 2025 Problem 2,ID=Jan25_2]
Let $n \geq 2$ be an integer, and $p \geq 3$ be a prime. Denote by $B$ the group of invertible $n \times n$ upper triangular matrices with entries in $\Z/p\Z$, and let $U$ be the subgroup of $B$ consisting of upper triangular matrices whose diagonal entries are all 1.
\begin{enumerate}
    \item What is the order of the quotient group $B/U$ as a function of $p$ and $n$?
    \item Give an example of a nontrivial element in the center of $U$.
    \item Prove that every nontrivial subgroup of $U$ has nontrivial center.
    \item Compute the center of $B$.
\end{enumerate}
\end{question}
\begin{solution}
\begin{enumerate}
    \item Any upper triangular matrix $t$ can be written as $t = du$ for a diagonal matrix $d$ and $u \in U$. Indeed, let $d$ be the diagonal matrix whose entries are the same as the diagonal entries of $t$, then $d\inv t \in U$. So, the elements of the quotient $B/U$ are represented by invertible diagonal matrices, of which there are $(p-1)^n$.
    \item Let $z$ be the matrix with 1s on the diagonal, a 1 in the top right corner, and 0s elsewhere. Then for any $u \in U$, $zu$ adds 1 to the top right entry of $u$, and so does $uz$, so $zu = uz$, so $z \in Z(U)$.
    \item The group $U$ has order a power of $p$, so every subgroup of $U$ also has order a power of $p$. Every $p$-group has nontrivial center, so every subgroup of $U$ must have nontrivial center.
    \item The center of $B$ consists of all scalar matrices. Clearly, the scalar matrices lie in $Z(B)$, so we need only show that any element of $Z(B)$ is a scalar matrix. First, any central element must be a diagonal matrix, for if $b$ is not a diagonal matrix, say $b_{ij} \neq 0$ for $i < j$. We can take the diagonal matrix $d = \underbrace{\operatorname{diag}(1,\dots,x,\dots 1)}_{\text{$x$ in position $i$}}$ for some $x \neq 1$ (which we can find since $p \geq 3$), and then the $(i,j)$ entry of $dbd\inv$ is $xb_{ij} \neq b_{ij}$, so $b \notin Z(B)$. Further, all the diagonal entries must be equal, for say that $b_{ii} \neq b_{jj}$. Let $e_{ij}$ be the matrix with 1s down the diagonal, a 1 in position $(i,j)$, and 0s elsewhere. Then $e_{ij} b \neq b e_{ij}$, so $b \notin Z(B)$. Hence, an element of $Z(B)$ must be a scalar matrix.
\end{enumerate}
\end{solution}

%
\begin{question}[subtitle={August 2024 Problem 4, August 2016 Problem 1},ID=Aug24_3+16_1]
Let $G$ be a group. By a maximal subgroup of $G$, we mean a subgroup $M \neq G$ such that the only subgroups containing $M$ are $M$ and $G$.
\begin{enumerate}
    \item Describe all the maximal subgroups of the dihedral group of order $2p$ where $p$ is an odd prime. How many are there?
    \item Show that if a finite group $G$ has only one maximal subgroup, then $G$ is cyclic.
    \item Show that if a maximal subgroup $M \triangleleft G$ is normal, then the index of $M$ in $G$ is finite and prime.
\end{enumerate}
\end{question}
\begin{solution}
\begin{enumerate}
    \item As in the previous problem, we have $D_p = \genset{r, f : r^p = f^2=rfrf=1}$ and every element of $D_p$ can be written as $r^i f^j$ where $0 \leq i \leq p-1$ and $0 \leq j \leq 1$. Since $\varphi(p)=p-1$, every rotation element $r^i$ with $1\leq i \leq p-1$ has order $p$, and each of these generates the same subgroup $\genset{r} \cong \Z/p\Z$. Every element of the form $r^i f$ with $0 \leq i \leq p-1$ has order $2$, and each generates its own subgroup $\genset{r^i f} \cong \Z/2\Z$. Since $\abs{D_p}=2p$, we know that any maximal subgroup must have order $2$ or $p$. Therefore, the maximal subgroups are exactly the single subgroup $\genset{r}$ of order $p$ and the $p$ subgroups $\genset{r^i f}$ of order $2$.
    \item Suppose $M < G$ is the only maximal subgroup of $G$. Then $M \neq G$, so there exists some $g \in G\setminus M$. Then $\genset{g}$ is a nontrivial subgroup of $G$, so it is either contained in a maximal subgroup or all of $G$ (since every proper subgroup must be contained in a maximal subgroup). Since, $g \notin M$, we must have that $\genset{g} = G$. Thus, $G$ is cyclic.
    \item Suppose $M \triangleleft G$ is maximal. Then the lattice theorem gives that the only subgroups of $G/M$ are the trivial group and the whole group. So, in particular $G/M$ has only one maximal subgroup, the trivial group, so $G/M$ is cyclic by the previous part. Let $n = \# G/M$, and notice that if $n$ were not prime then $G/M$ would have a subgroup of order $d\neq 1, n$ for some $d \mid n$. So, we conclude that in fact $n$ is prime, and since $\# G/M = [G:M]$ we conclude that the index of $M$ is also prime.
\end{enumerate}
\end{solution}

%
\begin{question}[subtitle = January 2024 Problem 3,ID=Jan24_3]
Let $G$ be a group acting transitively on a set $X$. Let $H \subset G$ be a normal subgroup. Then $H$ naturally acts on $X$, but the action need not be transitive. Prove that all orbits of $H$ on $X$ have the same cardinality. (You may assume $G$ and $X$ are finite if you like, but this is not necessary!)
\end{question}
\begin{solution}
Let's find a bijection between any two orbits. Let $x, y \in X$ be arbitrary, and denote their orbits under the action of $H$ by $Hx, Hy$, respectively. Since $G$ acts transitively, we know that there is some $g \in G$ so that $gx = y$, and hence $Hy = H(gx) = (Hg)x$, where by $(Hg)x$ I mean the set $\set{x' \in X: x' = hgx \text{ for some } h \in H}$. (The set $Hg$ is only a right coset, and yet it still makes sense for it to act on $X$.) Since $H$ is normal in $G$, the right coset $Hg$ is equal to the left coset $gH$, and thus $Hy = gHx$.

This gives us an easy bijection between the two: we get a map $Hx \to Hy$ given by $x' \mapsto gx'$, and a map $Hy \to Hx$ given by $x' \mapsto g\inv x'$. Our earlier argumentation shows that these maps indeed have the codomains we claim they have, and the two maps are clearly inverses of one another, which is enough to get what we need.
\end{solution}

%
\begin{question}[subtitle = {August 2023 Problem 5, edited},ID=Aug23_5]
By $S_n$ we mean the symmetric group on $n$ elements.
\begin{enumerate}
    \item Let $G$ be a subgroup of $S_n$. Suppose the action of $G$ on $\set{1,\dots,n}$ is transitive. (\textbf{NB} such a subgroup is called a ``transitive'' subgroup.) Prove that for any $x,y \in \set{1,\dots,n}$, the stabilizer of $x$ in $G$ and the stabilizer of $y$ in $G$ are conjugate subgroups of $G$.
    \item Prove that if $G$ is a transitive abelian subgroup of $S_n$, then $\abs{G} = n$.
    \item Give an example of an integer $n$ and a subgroup $G$ of $S_n$ such that $G$ is transitive and abelian but not cyclic.
    \item For which integers $n$ does $S_n$ contain a transitive abelian subgroup $G$ that is not cyclic? (Justify your answer.)
\end{enumerate}
\end{question}
\begin{solution}
\begin{enumerate}
    \item Since the action of $G$ is transitive, in particular there is some $g \in G$ so that $g \cdot x = y$. Then the conjugation functions
    \begin{align*}
        \phi\colon \Stab(x) &\to \Stab(y) & \psi\colon \Stab(y) &\to \Stab(x)\\
        h &\mapsto g h g\inv & f &\mapsto g\inv f g
    \end{align*}
    are in fact group homomorphisms, and they are clearly mutually inverse, so the two subgroups are conjugate.
    \item By part (a) we know that for any $x,y \in \set{1,\dots,n}$, $\Stab(x)$ is conjugate to $\Stab(y)$. If $G$ is abelian, then then any subgroup of $G$ is fixed by conjugation, so in fact if some $g \in G$ fixes some $x \in \set{1,\dots,n}$, it fixes every $x$, and therefore is the identity element. Therefore, we have
    \begin{align*}
        n &= \abs{G \cdot 1} \tag{since $G$ is transitive}\\
        &= [G:\Stab(1)] \tag{by Orbit-Stabilizer}\\
        &= \abs{G}. \tag{since $\Stab(1) = \set{e}$}
    \end{align*}
    \item Inside $S_4$, we have the abelian subgroup $\set{e,(12)(34),(13)(24),(14)(23)} \cong \Z/2\Z \times \Z/2\Z$.
    \item This occurs iff $n$ has a square factor. If $n$ has no square factor, then by the classification theorem there is only one abelian group of order $n$, and it is cyclic. On the other hand, if $n$ has a square factor, then there is some non-cyclic abelian group of order $n$, and by Cayley's theorem this group embeds as a transitive subgroup of $S_n$. (It acts transitively on itself by left multiplication.)
\end{enumerate}
\end{solution}

%
\begin{question}[subtitle = January 2023 Problem 2,ID=Jan23_2]
We say a group $G$ has \emph{big finite quotients} if for every $k > 0$, there is a normal subgroup $N$ of $G$ such that $G/N$ is finite and has order at least $k$.
\begin{enumerate}
    \item Give (with proof) an example of an infinite group that does \emph{not} have big finite quotients.
    \item Show that $\SL_2(\Z)$ (the group of $2 \times 2$ integer matrices with determinant 1) has big finite quotients.
\end{enumerate}
\end{question}
\begin{solution}
\begin{enumerate}
    \item The first infinite groups I think of are the additive groups of common rings. The group $\Z$ actually does have big finite quotients, so let's try $\Q$. The group $\Q$ is divisible, so all of its quotients are divisible, so all of its quotients are either trivial or infinite. Hence, $\Q$ does not have big finite quotients. %(Another possible solution is to use problem \ref{qu:PSL_simple}.)
    \item For every $p \in \Z$, there is a quotient map $\SL_2(\Z) \to \SL_2(\Z/p\Z)$ given by reducing each entry in the matrix mod $p$. One can compute that $\abs{\SL_2(\F_p)} = (p^2-1)(p^2-p)/(p-1) = p(p^2-1)$, so for any $k$ we can certainly just choose a $p > k$ and get a big quotient.
\end{enumerate}
\end{solution}

%LGFF
\begin{question}[subtitle = August 2022 Problem 2,ID=Aug22_2]
Let $\F_p$ be the finite field with $p$ elements; here $p$ is a prime number. Let $V = \F_p^2$, and recall that $G = \GL_2(\F_p)$ is the group of invertible linear transformations on $V$. $G$ acts on $V$ in the usual way (by multiplication).
\begin{enumerate}
    \item Describe the orbits of this action.
    \item Describe the stabilizer in $G$ of the vector
    \[\begin{pmatrix} 1 \\ 0 \end{pmatrix} \in V.\]
    \item Consider now the action of $G$ on $V \times V$ (acting independently on each of the two vectors). How many orbits does this action have?
    \item What is the cardinality of $G$? Remember to justify your answer.
\end{enumerate}
\end{question}
\begin{solution}
\begin{enumerate}
    \item There are two orbits, the zero vector, and everything else. For any two nonzero vectors $u_1,v_1$, there is some invertible matrix $A$ such that $Au_1 = v_1$; for example, we can choose complementary vectors $u_2,v_2$ so that $\set{u_1,u_2}$ and $\set{v_1,v_2}$ are both bases, and define an invertible linear map by $u_1 \mapsto v_1, u_2 \mapsto v_2$.
    \item The stabilizer of $\begin{pmatrix} 1\\0\end{pmatrix}$ consists of all matrices of the form
    \[\begin{pmatrix} 1 & a\\0 & b\end{pmatrix},\]
    where $a \in \F_p, b \in \F_p\mult$. Certainly, every matrix in the stabilizer must have $\begin{pmatrix} 1\\0\end{pmatrix}$ as the first column, and it is easy to verify that every such matrix is in the stabilizer.
    \item There's definitely an orbit corresponding to when $v_1 = 0$ and $v_2 = 0$. By our answer to part (a), there's also an orbit for when $v_1 \neq 0, v_2 = 0$, and another one for $v_1 = 0, v_2 \neq 0$. Indeed, somewhat more generally, if $v_1$ and $v_2$ are linearly dependent, say they are nonzero and $v_1 = c v_2$ with $c \in \F_p\mult$, then $A$ must take them to a new pair of vectors satisfying the same linear relation: $Av_1 = cAv_2$. This gives us $p-1$ more orbits. Finally, we have an orbit corresponding to when $v_1$ and $v_2$ are linearly independent vectors in $V$, because for any two such pairs there is a change-of-basis matrix taking one to the other. In total this gives us $p+3$ orbits.
    \item $(p^2-1)(p^2-p)$. One way to calculate this would be to use the calculation from part (b), together with the orbit-stabilizer theorem:
    \begin{align*}
        \abs{G} &= \abs{\Stab \begin{pmatrix} 1\\0\end{pmatrix}} \cdot \abs{G \cdot \begin{pmatrix} 1\\0\end{pmatrix}}\\
        &= p(p-1) \cdot (p^2-1).
    \end{align*}
\end{enumerate}
\end{solution}

%
\begin{question}[subtitle = January 2022 Problem 1,ID=Jan22_1]
On this problem, only the answer will be graded. Consider the symmetric group $S_4$.
\begin{enumerate}
    \item What is the order of $S_4$?
    \item What is the order of its center $Z(S_4)$?
    \item What is the order of its commutator subgroup $[S_4, S_4]$?
    \item Is $S_4$ a simple group?
    \item Is $S_4$ a solvable group?
    \item How many conjugacy classes does $S_4$ have?
\end{enumerate}
\end{question}
\begin{solution}
\begin{enumerate}
    \item 24
    \item 1
    \item 12 ($[S_4,S_4]=A_4$; this is always the case)
    \item No, $A_4$ is a nontrivial normal subgroup.
    \item Yes, since $A_4$ contains a normal subgroup isomorphic to $\Z/2\Z \times \Z/2\Z$, and the quotient is isomorphic to $\Z/3\Z$.
    \item 5, one for each partition of 4.
\end{enumerate}
\end{solution}

%
\begin{question}[subtitle = August 2021 Problem 3,ID=Aug21_3]
Let $G$ be a finite group. For given $n >1$, we say that $G$ \emph{admits roots of degree $n$} if, for every $x\in G$, there exists $y\in G$ such that $y^n=x$.
\begin{enumerate}
    \item Show that if $n$ is coprime to $\abs{G}$, $G$ admits roots of degree $n$.
    \item Conversely, prove that if $G$ admits roots of degree $n$, then $n$ is coprime to $\abs{G}$.
\end{enumerate}
\end{question}
\begin{solution}
\begin{enumerate}
    \item We will prove the proposition in the case where $G$ is cyclic, and the result will follow because each $x \in G$ lives in the cyclic subgroup $\genset{x}$. So, let $G$ be cyclic and generated by $x$, and consider the group homomorphism $\phi_n\colon G \to G$ given by $x \mapsto x^n$. We want to show that this map is surjective. Since $G$ is finite, it is enough to show that this map is injective, and hence it is enough to show that $\ker \phi_n = \set{e}$. Suppose $x \in \ker \phi_n$, then $x^n = e$, so the order $x$, $\ord x$, divides $n$. But also $\ord x$ divides $\abs{G}$. Since $n$ is coprime to $\abs{G}$ we must have that $\ord x = 1$, i.e.\ $x = e$. Thus, $\ker \phi_n$ is trivial.
    \item We will prove this by contradiction: suppose $\gcd(n,\abs{G}) > 1$ and $G$ admits roots of degree $n$. If $G$ admits roots of degree $n$, it also admits roots of degree $p$ for every $p \mid n$, so it suffices to consider the case where $n = p$ for some prime $p$ dividing $\abs{G}$. Let $p \mid \abs{G}$, and let $g \in G$ generate a cyclic subgroup of order $p^k$ which is not contained in any cyclic subgroup of order $p^{k+1}$. By Cauchy's theorem, there a cyclic subgroup of order $p$, and we are just taking a maximal cyclic $p$-group containing such a subgroup. Since $G$ admits roots of degree $p$, we can find $h \in G$ with $h^p = g$. First, notice that for any $k$, $(g^k)^p$ lies in a subgroup of index $\geq p$ inside $\genset{g}$, and therefore $h \notin \genset{g}$. But, on the other hand, $\ord h = p^{k+1}$, and by construction $g$ is not contained in any cyclic group of order $p^{k+1}$, a contradiction. Therefore, $G$ cannot admit roots of degree $p$.
\end{enumerate}
\end{solution}

%
\begin{question}[subtitle = January 2020 Problem 5,ID=Jan20_5]
\begin{enumerate}
    \item Find a nontrivial finite subgroup of $\GL_2(\R)$.
    \item Find an infinite, abelian, but non-cyclic subgroup of $\GL_2(\R)$.
    \item Find a proper, normal subgroup of $\GL_2(\Z)$ of finite index.
\end{enumerate}
\end{question}
\begin{solution}
\begin{enumerate}
    \item Consider the subgroup of order 2 $\set{I,-I}$.
    \item Consider the subgroup consisting of diagonal matrices.
    \item For each prime $p$ there is a map $\GL_2(\Z) \to \GL_2(\F_p)$ given by reducing each entry mod $p$. The kernel of such a map must be proper and normal, and since the target is a finite group it must be of finite index.
\end{enumerate}
\end{solution}

%LGFF
\begin{question}[subtitle = August 2019 Problem 5,ID=Aug19_5]
Let $G$ be the group of $2 \times 2$ invertible upper triangular matrices over the field $\F_p$.
\begin{enumerate}
    \item Prove that $G$ has only one subgroup of order $p$, and that this subgroup is normal. (Hint: first show that an element of order $p$ must have 1's on the diagonal.)
    \item Prove that $G$ is solvable by exhibiting a homomorphism $f$ from $G$ to an abelian group $A$ such that the kernel of $f$ is also abelian.
    \item Prove that if $p \neq 2$, $G$ is not nilpotent.
\end{enumerate}
\end{question}
\begin{solution}
\begin{enumerate}
    \item Given a matrix
    \[A = \begin{pmatrix}a&b\\0&c \end{pmatrix},\]
    we have that the diagonal entries of $A^n$ are $a^n$ and $c^n$. Thus, if $A$ has order $p$, then $a^p = c^p = 1$, but by Fermat's Little Theorem, $a^p = a$, and so $a = c = 1$. Thus, the $p$ matrices of the form
    \[\begin{pmatrix} 1 & b\\ 0 & 1 \end{pmatrix}\]
    form the only subgroup of order $p$. Since conjugating this subgroup would take it to another subgroup of order $p$, and we have determined it is the unique subgroup of order $p$, we conclude that this subgroup must be fixed by conjugation, hence normal.
    \item Let the subgroup from part (a) be $U$, and let $T$ be the subgroup of diagonal matrices. Then $U \cap T = \set{I}$, and we can write any matrix
    \[\begin{pmatrix} a&b\\0&c \end{pmatrix} = \begin{pmatrix} a& 0\\0 & c\end{pmatrix} \begin{pmatrix} 1 & a\inv b\\0& 1 \end{pmatrix}\]
    where the first factor is in $T$ and the second is in $U$. Thus, $G = U \rtimes T$, and so $G/U$ is isomorphic to $T$, which is abelian. Thus, $f\colon G \to G/U$ is the homomorphism we want.
    \item Assume $p \neq 2$. This is probably easiest to show using the lower central series. Notice that the diagonal entries of any commutator $ABA\inv B\inv$ are going to be 1, so actually every commutator falls inside $U$. So, $[G,G]$ is a subgroup of $U$, and it isn't the trivial subgroup since $G$ is not abelian, so the only other option is $[G,G] = U$. But now, the second term of the lower central series $G_2 = [G,U]$ is still equal to $U$, since for example the matrices
    \[A = \begin{pmatrix} 1&1\\0&1\end{pmatrix}, \qquad B = \begin{pmatrix}2&0\\0&1 \end{pmatrix}\]
    do not commute (recall that we've assumed $p \neq 2$). Thus, $ABA\inv B\inv$ is a nontrivial commutator in $[G,U]$, so $[G,U]$ is a nontrivial subgroup of $U$, hence is all of $U$.
\end{enumerate}
\end{solution}

%Pretty sure I can improve the solution to part (d).
\begin{question}[subtitle = January 2019 Problem 2,ID=Jan19_2]
If $G$ is a group and $x, y\in G$, then we let $\genset{x, y}$ denote the subgroup of $G$ generated by $x$ and $y$. For $x \in G$, define $\Cyc_G(x) \coloneqq\set{y\in G\ :\ \genset{x, y} \text{ is cyclic}}$.
\begin{enumerate}
    \item Show by example that $\Cyc_G(x)$ need not be a subgroup of $G$.
    \item Let $\Cyc(G) =\bigcap_{x\in G} \Cyc_G(x)$. Show that $\Cyc(G)$ is a subgroup of $G$.
    \item Show that $\Cyc(G)$ lies in the center of $G$.
    \item Assume now that $G$ is finite. Show that $\Cyc(G)$ is cyclic.
\end{enumerate}
\end{question}
\begin{solution}
\begin{enumerate}
    \item When thinking of a counterexample, it's helpful to rephrase the definition of $\Cyc_g(x)$. For any cyclic subgroup $\genset{g}$ which contains $x$, we will have that $\genset{x,y} \leq \genset{g}$ for all $y \in \genset{g}$, so $\genset{x,y}$ will also be cyclic. Thus, every $y \in \genset{g}$ is in $\Cyc_g(x)$, and this holds for all $\genset{g}$, so $\Cyc_g(x)$ is the union of all the cyclic subgroups containing $x$:
    \[\Cyc_g(x) = \bigcup_{\substack{g \in G\\ x = g^k}} \genset{g}.\]
    However, a union of subgroups need not be a subgroup, so we should try to think of examples where that fails.
    
    As a specific example, $G = \Z/2\Z \times \Z$ and $x = (0,2)$. Then $(1,1)$ and $(0,1)$ are in $\Cyc_G(x)$ since $\genset{(0,2),(1,1)} = \genset{(1,1)})$ and $\genset{(0,2), (0,1)} = \genset{(0,1)}$, but $(1,0)=(1,1)-(0,1)$ is not in $\Cyc_G(x)$ since $\genset{(0,2), (1,0)} = \Z/2\Z \times 2\Z$.
    \item It's helpful to think about what it means to say ``$y$ is an element of $\Cyc(G)$''. If $y \in \Cyc(G)$, then $y \in \Cyc_G(x)$ for every $x \in G$, so $\genset{x,y}$ is cyclic for every $x \in G$. (And the implications go the other way as well.)
    
    We first show that $\Cyc(G)$ is closed under taking inverses. Let $y \in \Cyc(G)$. We need to show that for each $x \in G$, $\genset{x, y^{-1}}$ is cyclic. But $\genset{x, y^{-1}} = \genset{x, y}$, which is cyclic by assumption, so $y^{-1} \in \Cyc(G)$.
    
    Now we show that $\Cyc(G)$ is closed under multiplication. It will be convenient to define a function (not a homomorphism) $g\colon G \times \Cyc(G) \to G$ by choosing $g_y(x)$ to be any generator for the cyclic group $\genset{x,y}$. (I put the $y$ in the subscript to make the following easier to read.) Now, let $y,z \in \Cyc(G)$, and let $x \in G$ be arbitrary. Observe that $\genset{x,y,z} = \genset{g_y(x),z} = \genset{g_z(g_y(x))}$ is cyclic. Since $y, z \in \genset{x,y,z}$, $yz \in \genset{x,y,z}$, so $\genset{x,yz} \leq \genset{x,y,z}$ is a subgroup of a cyclic group, hence is itself cyclic. Thus, $yz \in \Cyc(G)$.
    \item Let $y \in \Cyc(G)$. We want to show that $y \in Z(G)$, that is, that $y$ commutes with every element of $G$. We know that for any $x \in G$, $\genset{x, y}$ is cyclic, and hence abelian. Therefore, $x$ and $y$ must commute.
    \item We will reuse our function $g\colon G \times \Cyc(G) \to G$ from part (b). Since $G$ is finite, we can find finitely many $y_1,y_2\dots,y_n$ that generate $\Cyc(G)$. But then, similar to part (b),
    \[\Cyc(G) = \genset{y_1,\dots,y_n} = \genset{g_{y_2}(y_1),y_3,\dots,y_n} = \dots = \genset{g_{y_n}\circ g_{y_{n-1}} \circ \dots g_{y_2}(y_1)}.\]
    Hence, $\Cyc(G)$ is cyclic.
    
    % First, observe that $\Cyc(G) \leq \Cyc(Z(G))$, since certainly if $\genset{x,y}$ is cyclic for all $x \in G$, then $\genset{z,y}$ is cyclic for all $z \in Z(G)$. Since a subgroup of a cyclic group is cyclic, it is enough to show that $\Cyc(Z(G))$ is cyclic. Since $Z(G)$ is abelian, it suffices to show that $\Cyc(G)$ is cyclic when $G$ is abelian. By the fundamental theorem of finitely generated abelian groups, we can write $G = \prod_{i=1}^k \Z/p_i^{e_i}\Z$ for some primes $p_i$ and exponents $e_i \geq 1$. Let $I$ be the set of indices $i$ such that $p_i \neq p_j$ for any index $j \neq i$. We will show that $\Cyc(G) = \prod_{i \in I} \Z/p_i^{e_i}\Z$, which is a cyclic group since no $p_i$ is repeated.
    
    % To show that $\Cyc(G) \subseteq \prod_{i \in I} \Z/p_i^{e_i}\Z$, suppose instead that there is a $y \in \Cyc(G)$ that is not in the RHS. Then, there must exist indices $i,j \in \set{1,\dots,k}\setminus I$ such that $p_i = p_j = p$ where $y$ has a nonzero component in at least one of the factors $\Z/p^{e_i}\Z$ or $\Z/p^{e_j}\Z$ of $G$. Let $(a,b)$ be the image of $y$ under the projection from $G$ to $\Z/p^{e_i}\Z \times \Z/p^{e_j}\Z$, and assume without loss of generality that $a \neq 0$. Let $x \in G$ be any element that projects to $(0,1)$. Since $\genset{x, y}$ is cyclic, its image must also be cyclic. But, $\genset{(0,1), (a,b)} \cong \Z/(p^{e_i}/\gcd(a,p^{e_i}))\Z \times \Z/p^{e_j}\Z$, which is not cyclic since $p^{e_i}/\gcd(a,p^{e_i})$ and $p^{e_j}$ are nonzero powers of $p$.
    
    % To show the reverse inclusion, let $A = \prod_{i \in I} \Z/p_i^{e_i}\Z$, and write $G = A \times B$, where $B$ consists of all of the remaining factors for $i \neq I$. Let $(a,0) \in A$. We want to show that $(a,0) \in \Cyc(G)$, i.e.\ that for any $x = (x_A, x_B) \in G$ the subgroup $\genset{(x_A,x_B), (a,0)}$ is cyclic. Consider that
    % \[\genset{(x_A,x_B), (a,0)} \leq \genset{(x_A,0),(0,x_B), (a,0)} = \genset{\genset{(x_A,0),(a,0)}, (0,x_B)}.\]
    % Since $\genset{(x_A,0),(a,0)}$ is a subgroup of $A$, it is cyclic, so let $\genset{(x_A,0),(a,0)} = \genset{a'}$ and suppose its order is $n$. Note that $n$ is divisible only by primes $p_i$ with $i \in I$. Then $(0,x_B)$ also generates a cyclic subgroup of $G$, and its order must be coprime to $n$. Therefore,
    % \[\genset{(x_A,x_B),(a,0)} \leq \genset{(a',0),(0,b)} \leq \genset{a'} \times \genset{b}.\]
    % Notice that the RHS is cyclic since it is a product of cyclic groups of coprime order. Thus, the LHS is cyclic as a subgroup of a cyclic group. Hence, $\Cyc(G) = A$.
\end{enumerate}
\end{solution}

%An old favorite of mine! The first practice qual problem I ever did!
\begin{question}[subtitle = August 2018 Problem 2,ID=Aug18_2]
For a finite group $G$, denote by $s(G)$ the number of subgroups of $G$.
\begin{enumerate}
    \item Show that $s(G)$ is finite.
    \item Show that if $H$ is a nontrivial normal subgroup of $G$, then $s(G/H) < s(G)$.
    \item Show that $s(G) = 2$ if and only if $G$ is cyclic of prime order.
    \item Show that $s(G) = 3$ if and only if $G$ is cyclic of order $p^2$ for a prime $p$.
\end{enumerate}
\end{question}
\begin{solution}
\begin{enumerate}
    \item Certainly every subgroup is also a subset, and there are at most $2^{\abs{G}}$ subsets of $G$, so $s(G) \leq 2^{\abs{G}} < \infty$.
    \item By the lattice isomorphism theorem, the subgroups of $G/H$ are in bijection with subgroups of $G$ containing $H$. Certainly the number of subgroups of $G$ containing $H$ is at most the number of subgroups of $G$, and since $H$ is nontrivial we know the trivial subgroup does not contain $H$, which makes the inequality strict.
    \item Certainly, if $G$ is cyclic of order $p$, then its only subgroups are $G,\set{e}$. On the other hand, suppose $s(G) = 2$. Since $s(G) > 1$, $G$ is not trivial. Let $p$ be a prime dividing $\abs{G}$, then by Cauchy's theorem there is a $g \in G$ of exact order $p$. But then $\genset{g}$ and $\set{e}$ are 2 subgroups of $G$, so we have already found all the subgroups of $G$.
    \item Certainly, if $G$ is cyclic of order $p^2$ generated by $g$, then its only subgroups are $G, \genset{g}, \set{e}$. On the other hand, suppose $s(G) = 3$, so that $G$ has exactly 1 nontrivial proper subgroup. As before, we can find an element $g$ of exact order $p$. Furthermore, $p$ is the only prime which divides $\abs{G}$, as if another prime $q$ divided $\abs{G}$ we would have another nontrivial proper subgroup, but we know $G$ has a unique such subgroup. Thus, $\abs{G} = p^n$ for some $n$. Now, let $h \in G \setminus \genset{g}$, then $h$ generates a nontrivial subgroup, so that subgroup cannot be proper, so $\genset{h} = G$, and $G$ is cyclic. Since $G$ is cyclic of order $p^n$, it has $n+1$ subgroups, generated by $h^{p^k}$ for $k \leq n$. Thus, $n+1 = 3$, so $n = 2$.
\end{enumerate}
\end{solution}

%I changed part (a) to ask a more straightforward question.
\begin{question}[subtitle = {January 2018 Problem 4, modified},ID=Jan18_4]
Let $G$ be a finite group. Denote by $\Aut(G)$ the group of automorphisms of $G$, and by $Z(G) \subset G$ the center of $G$.
\begin{enumerate}
    \item Let $\Inn(G) \subset \Aut(G)$ be the subgroup consisting of automorphisms coming from conjugation, i.e.\ automorphisms of the form $x \mapsto gxg\inv$ for some $g \in G$. Prove that $\Inn(G)$ is isomorphic to $G/Z(G)$.
    \item Show that if $G/Z(G)$ is cyclic, then $G$ is abelian.
    \item Show that if $\Aut(G)$ is cyclic, then $G$ is abelian.
    \item Show that if $G$ is abelian, then $x \mapsto x\inv$ is an automorphism of $G$. (\textbf{NB} in fact this is an if and only if.)
    \item Deduce that there is no group $G$ such that $\Aut(G)$ is a nontrivial cyclic group of odd order.
\end{enumerate}
%I remember I looked a little bit into what can be said about when Aut(G) is abelian. I think I found something interesting, but I no longer remember what it was...
\end{question}
\begin{solution}
\begin{enumerate}
    \item Denote by $\conjg{g}$ the ``conjugation by $g$'' automorphism, $\conjg{g} \colon G \to G$, $\conjg{g}(x) = gxg\inv$. Define a map $G \to \Inn(G)$ by $g \mapsto \conjg{g}$. We will show that the kernel of this map is equal to $Z(G)$. Certainly, if $z \in Z(G)$, then $zxz\inv = x$ for all $x$, so $\conjg{z} = \text{id}$. On the other hand, suppose $\conjg{g} = \text{id}$, that is, that $gxg\inv = x$ for all $x$. Then multiplying on the right by $g$ gives us $gx = xg$ for all $x$, so $g \in Z(G)$.
    \item Suppose $G/Z(G)$ is cyclic with generator $gZ(G)$. Then any $x,y \in G$ can be written uniquely as $x = g^{n_1} z_1$, $y = g^{n_2} z_2$, for $z_1,z_2 \in Z(G)$. But then
    \[xy = g^{n_1} z_1 g^{n_2} z_2 = g^{n_1} g^{n_2} z_1 z_2 = g^{n_2} g^{n_1} z_2 z_1 = g^{n_2} z_2 g^{n_1} z_1 = yx,\]
    and $G$ is abelian.
    \item If $\Aut(G)$ is cyclic, then its subgroup $\Inn(G) \cong G/Z(G)$ is also cyclic.
    \item Notice that $(xy)\inv = y\inv x\inv = x\inv y\inv$ since $G$ is abelian.
    \item By the proceeding, if $\Aut(G)$ is nontrivial and cyclic, then $G$ is abelian. Then $x \to x\inv$ is an element of $\Aut(G)$ of order dividing 2, so either $2 \mid \abs{\Aut(G)}$ or every element of $G$ has order exactly 2. In the other case, we will have $G \cong (\Z/2\Z)^n$ for some $n>1$ by the classification of finite abelian groups. But then there is an automorphism defined by swapping $e_1$ and $e_2$ and leaving the other basis elements fixed, and this is an automorphism of order 2, so even in this case $2 \mid \abs{\Aut(G)}$. Thus, it is impossible for $\Aut(G)$ to be a nontrivial cyclic group of odd order. 
\end{enumerate}
\end{solution}

%This is the most recent qual problem that mentions a Sylow subgroup. I think August 2019 Problem 5 is the only more recent problem where I think knowing anything about the Sylow theorems being helpful, and that's not much more recent.
\begin{question}[subtitle = August 2017 Problem 3,ID=Aug17_3]
What is the smallest $n$ such that the 3-Sylow subgroup of $S_n$ is non-abelian? (You may use the Sylow theorem that all Sylow subgroups are conjugate, so that one 3-Sylow subgroup is non-abelian if and only if they all are.)
\end{question}
\begin{solution}
Any group of order $p$ or $p^2$ is abelian, so the smallest $n$ such that the 3-Sylow subgroup of $S_n$ could be nonabelian is $n=9$. We have to be a little careful showing that in this case the 3-Sylow subgroup is nonabelian. We can certainly find two elements whose orders are powers of 3 and don't commute, but there's always the possibility that they live in different Sylow subgroups, in which case we will not have demonstrated that the 3-Sylow subgroup is nonabelian. (As an aside: by other considerations, there are at least 4 3-Sylow subgroups, so if we were to choose our elements randomly we would get two elements in the same Sylow subgroup at best 1/4 of the time. Those aren't great odds.)

However, we can say that a 3-Sylow subgroup must contain a 3-cycle $(abc)$, because certainly every 3-cycle lives in \emph{some} 3-Sylow subgroup, but all such subgroups are conjugate, and conjugation does not change cycle type, so each 3-Sylow subgroup gets at least one of them. For the same reason, we can see that every 3-Sylow subgroup has a 9-cycle. However, a 3-cycle never commutes with a 9-cycle. Indeed, given a 3-cycle $\tau$ and a 9-cycle $\nu$, there is always some $n \in \set{1,\dots,9}$ so that $\tau(n) \neq n$ and $\tau\circ \sigma(n) = \sigma(n)$. (E.g. if $\tau = (123)$ and $\nu = (123456789)$, then $n = 3$ is such a number, and is the only such number.) But then clearly $\sigma\circ \tau(n) \neq \tau\circ \sigma(n)$, so these two do not commute.
\end{solution}

%
\begin{question}[subtitle = January 2017 Problem 1,ID=Jan17_1]
Give an example of each of the following.
\begin{enumerate}
    \item A group $G$ with a normal subgroup $N$ such that $G$ is not a semidirect product $N \rtimes G/N$.
    \item A finite group $G$ that is nilpotent but not abelian.
    \item A group $G$ whose commutator subgroup $[G,G]$ is equal to $G$.
    \item A non-cyclic group $G$ such that every Sylow subgroup of $G$ is cyclic.
    \item A transitive action of $S_3$ on a set $X$ of cardinality greater than 3.
\end{enumerate}
\end{question}
\begin{solution}
\begin{enumerate}
    \item The smallest such example is $G = \Z/4\Z$, $N = 2\Z/4\Z$. Then $G/N \cong \Z/2\Z$, but there is only one nontrivial map $\Z/2\Z \to \Z/4\Z$, and it is not a splitting for the quotient map. Therefore, $\Z/4\Z$ is not a semidirect product of $\Z/2\Z$ and $\Z/2\Z$. A non-abelian example would be the quaternion group $Q$ with subgroup $N = \set{\pm 1}$. An infinite example would be $\Z$ with subgroup $N = 2\Z$.
    \item The smallest examples are $D_4$ and $Q$, the quaternion group.
    \item The smallest such example is $G = A_5$, which is simple and therefore its commutator subgroup must be equal to the whole group.
    \item Take $G = D_n$ for $n$ odd.
    \item Consider the action of $S_3$ on itself by left multiplication.
\end{enumerate}
\end{solution}

%
\begin{question}[subtitle=January 2016 Problem 4,ID=Jan16_4]
Let $D_k$ be the dihedral group of order $2k$ where $k\geq 3$.
\begin{enumerate}
    \item Show that the number of automorphisms of $D_k$ is equal to $k\cdot \varphi(k)$, where $\varphi$ is Euler's $\varphi$-function.
    \item Let $\Aut(D_k)$ denote the group of automorphisms of $D_k$. What is the structure of $\Aut(D_k)$? Describe the group as explicitly as you can.
\end{enumerate}
\end{question}
\begin{solution}
\begin{enumerate}
    \item Recall that $D_k$ has presentation given by
    \[\genset{r, f : r^k = f^2 = rfrf = 1}\]
    where $r$ is a rotation of the $k$-gon that sends each vertex to an adjacent one and $f$ is a reflection of the $k$-gon through a line of symmetry. This presentation shows that every element of $D_k$ can be written in the form $r^i f^j$, where $0 \leq i \leq k-1$ and $0 \leq j \leq 1$. Notice that the only elements of $D_k$ of order $k$ are those of the form $r^i$ where $i$ is relatively prime to $k$, and the only elements of order $2$ are those of the form $r^i f$.
    
    Any element of $\Aut(D_k)$ is determined by where it sends $r$ and $f$, and $r$ must be sent to another element of order $k$ and $f$ must be sent to another element of order $2$.Therefore, there are $\varphi(k)$ choices for the image of $r$ and $k$ choices for the image of $f$. This shows that $\abs{\Aut(D_k)} \leq k\cdot \varphi(k)$. It remains to check that each of these choices gives an automorphism of $D_k$. It is routine to check that each map is a homomorphism (just need to check the relations still hold), and given the map $\varphi$ where $\varphi(r) = r^i$ and $\varphi(f) = r^j f$, we can construct an inverse. Since $\gcd(i,k)=1$, there exist $a,b \in \Z$ such that $ai+bk = 1$, and then $\varphi^{-1}$ is given by $r \mapsto r^a$ and $f \mapsto r^{-ja}f$. Thus, $\abs{\Aut(D_k)}= k\cdot \varphi(k)$.
    \item By part (a), we know that any automorphism of $D_k$ is determined by a pair of integers $(j,i)$ where $0 \leq j \leq k-1$ and $\gcd(i,k)=1$. Let's next understand the multiplication in this group. Given two automorphism determined by $(j,i)$ and $(n,m)$, we see that their composition maps $r \mapsto r^m \mapsto r^{im}$ and $f \mapsto r^n f \mapsto r^{in+j}f$. Therefore, we have $(j,i) \circ (n,m) = (in+j, im)$. This looks like it might indicate a semi-direct product. Indeed, if we consider the map $\alpha: (\Z/k\Z)^* \to \Aut(\Z/k\Z)$ which sends $i$ to the multiplication by $i$ map, we see that $\Aut(D_k) \cong \Z/k\Z \rtimes (\Z/k\Z)^*$ with respect to $\alpha$.
\end{enumerate}
\end{solution}